% Chapter 4 of Singular Learning Theory for the Rest of Us: the volume of nearly true parameters.
\sltchapter{Volume}{The learning coefficient as a scaling exponent}{$y=a^2x_1+b^2x_2$}{ch:volume}
\newtonfoot{0/4,4/0}{2}{\tfrac12}{a^4+b^4}
\begin{refsection}[chapters/04-volume.bib]

\begin{chaptersummary}
How much prior mass do the parameters hold that are almost as good as the truth, and how fast does it shrink? For a regular model it falls like $\eps^{d/2}$; for the quartic cup $y=a^2x_1+b^2x_2$, whose only true parameter has a zero Hessian, we compute $V(\eps)=1.311\,\eps^{1/2}$ exactly, and the volumes of the other running models by scaling and by logarithmic coordinates. We prove that $V(\eps)\sim C\eps^\lambda(\log\frac1\eps)^{m-1}$ gives the population free energy $\lambda\log n-(m-1)\log\log n+O(1)$, that a normal-crossing function has $\lambda=\min_j\frac{h_j+1}{2k_j}$, and that $-\lambda$ is the largest pole of the zeta function and $m$ its order, in both directions. Monte Carlo with up to $10^8$ prior draws recovers every exponent to within $0.001$ once the logarithm is modeled, and shows the price of volumes: about $100/V(\eps)$ draws for $10\%$ accuracy, which no model with many parameters can pay.
\end{chaptersummary}

\section{The question: how much room near the truth?}\label{sec:volume:question}

A trained network sits in a region of weight space where the loss is nearly minimal, and deep learning has long asked how large such regions are. \Textcite{HochreiterSchmidhuber1997} searched for \term{flat minima}\index{flat minima}, large connected regions in which the error stays approximately constant; later measures of flatness mostly use the eigenvalues of the Hessian, which \textcite{DinhEtAl2017} showed to depend on the parametrization.\sidenote{A prior mass does not: if $w=g(u)$ is a smooth bijection and the prior is carried along as a density, the set $\{K<\eps\}$ and its prior mass are the same in both coordinates.} A Bayesian asks the question with a prior. If $K(w)=L(w)-L(w_0)$ is the excess loss of a parameter over the truth and $\varphi$ the prior, the parameters that are almost as good as the truth hold the prior mass
\begin{equation}\label{eq:volume:V}
V(\eps)=\int_{\{K(w)<\eps\}}\varphi(w)\dd w,\qquad\eps>0,
\end{equation}
the \term{volume function}\index{volume function} of $K$ with respect to $\varphi$.\sidenote{For the uniform prior on a box, $V(\eps)$ is the Lebesgue volume of $\{K<\eps\}$ divided by the volume of the box. Watanabe, and \textcite{LauEtAl2025}, speak of volumes, and so do we.}

\Cref{ch:coefficient} met this function in one dimension: when the likelihood depends on $w$ only through a scalar coefficient $c$, $V(\eps)$ is the prior mass of the window $\abs{c-c_0}<\sqrt{2\eps}$, and the free energy is, up to a bounded term, minus the logarithm of that mass at $\eps=1/n$ (\cref{thm:coefficient:free,prop:coefficient:zeta}). The model of this chapter, the \term{quartic cup}\index{quartic cup} $y=a^2x_1+b^2x_2+\varepsilon$ with $x\sim\Normal(0,I_2)$, unit noise and the uniform prior on $[-1,1]^2$, fitted to pure noise, has two coefficients and the loss
\[
K(a,b)=\tfrac12\E_x\bigl[(a^2x_1+b^2x_2)^2\bigr]=\tfrac12\bigl(a^4+b^4\bigr).
\]
Its single true parameter $(0,0)$ has a zero Hessian, so the pair is singular. The free energy is no longer an integral in one variable, but the population free energy $-\log\int e^{-nK}\varphi$ depends on $K$ only through its law under the prior, which $V$ encodes: the volume function replaces the coefficient density.

\begin{regularanswer}[the volume near a nondegenerate minimum]
Let $K$ be continuous on a compact $W$ with a single zero $w_0$, an interior point near which $K$ is three times continuously differentiable with a positive definite Hessian $J$, and let $\varphi$ be continuous and positive at $w_0$. Then $\{K<\eps\}$ is, for small $\eps$, nearly the ellipsoid $\tfrac12(w-w_0)^\top J(w-w_0)<\eps$, and
\[
V(\eps)\sim\varphi(w_0)\,\frac{\pi^{d/2}\,(2\eps)^{d/2}}{\Gamma(\frac d2+1)\sqrt{\det J}}\qquad(\eps\to0):
\]
half a power of $\eps$ per parameter. For the regular plane $y=ax_1+bx_2$ and the uniform prior on $[-1,1]^2$, $V(\eps)=\frac\pi2\eps$ exactly for $\eps\le\frac12$.
\end{regularanswer}

The constant is the volume of an ellipsoid (\cref{def:toolkit:gamma}); its logarithm contains the curvature term $-\frac12\log\det J$ of Laplace's formula (\cref{thm:regular:laplace}). For the quartic cup the formula fails at its first step: $J=0$, and the formula returns an infinite volume, although $V\le1$. For the product model $y=abx$ the true parameters form the cross $ab=0$, on whose arms $J$ has rank one and at whose centre it vanishes. The question of this chapter is: \emph{how fast does $V(\eps)$ shrink for a singular model, what does the rate control, and how is it related to the zeta function, through which the book defines the learning coefficient?}

\section{By hand: four volumes}\label{sec:volume:byhand}

\begin{computation}{The volume of the quartic cup}\label{comp:volume:cup}
\step \emph{The set is a scaled copy of one set.} Since $K(ta,tb)=t^4K(a,b)$, the set $\{K<\eps\}$ is $\{u^4+v^4<1\}$ scaled by $(2\eps)^{1/4}$ in both directions, and for $\eps\le\frac12$ it lies in the square.
\step \emph{That set has an explicit area.} By symmetry and the substitution $s=v^4$, the area of $\{u^4+v^4<1\}$ is
\[
A_4=4\int_0^1(1-v^4)^{1/4}\dd v=\int_0^1s^{-3/4}(1-s)^{1/4}\dd s=B\bigl(\tfrac14,\tfrac54\bigr),
\]
which is $\Gamma(\frac14)^2/(2\sqrt\pi)=3.7081$ (\cref{def:toolkit:gamma}).
\step \emph{The volume is an exact power.} Scaling by $(2\eps)^{1/4}$ multiplies areas by $(2\eps)^{1/2}$, and the prior density is $\frac14$, so for $0<\eps\le\frac12$
\[
V(\eps)=\tfrac14A_4\,(2\eps)^{1/2}=C\eps^{1/2},\qquad C=\frac{\Gamma(\frac14)^2}{4\sqrt{2\pi}}=1.3110 .
\]
Dividing $\eps$ by ten divides the volume by $\sqrt{10}$, the regular plane's by ten.
\step \emph{The state density and the zeta function.} The law of $K$ under the prior has the density $v(t)=V'(t)=\frac C2t^{-1/2}$ on $(0,\frac12)$, and the zeta function $\zeta(z)=\int K^z\varphi$ integrates $t^z$ against this law:
\[
\zeta(z)=\frac C2\int_0^{1/2}t^{z-1/2}\dd t+\int_{\{K\ge1/2\}}K^z\varphi=\frac{C\,2^{-z-1/2}}{2z+1}+\int_{\{K\ge1/2\}}K^z\varphi .
\]
Since $\frac12\le K\le1$ on $\{K\ge\frac12\}$, the last integral is an entire function of $z$. So $\zeta$ has one pole, simple, at $-\frac12$, and $(z+\frac12)\zeta(z)\to\frac C2=0.6555$.
\step \emph{The population free energy.} The partition function factorizes, $Z(n)=\int e^{-nK}\varphi=\bigl(\int_0^1e^{-na^4/2}\dd a\bigr)^2$; extending the integral to $(0,\infty)$ changes it by a relative error below $e^{-n/2}$ for $n\ge3$, and $\int_0^\infty e^{-na^4/2}\dd a=\Gamma(\frac54)(2/n)^{1/4}$. So, up to an exponentially small error,
\[
\bar F(n)=-\log Z(n)=\tfrac12\log n-\log\bigl(\sqrt2\,\Gamma(\tfrac54)^2\bigr)=\tfrac12\log n-0.1500 .
\]
At $n=10^4$ the cup pays $4.455$ nats; the regular plane, with $\bar F(n)=\log n-\log\frac\pi2$, pays $8.759$.
\step \emph{The same volume, in the coefficients.} In the coefficients $c=(a^2,b^2)$ the loss is $\frac12\abs c^2$, the loss of the regular plane; what differs is the prior. Each $c_j=a_j^2$ has the density $\frac12c_j^{-1/2}$ on $(0,1]$, so $c$ has the density $(4\sqrt{c_1c_2})^{-1}$, infinite at the truth. Over the quarter disc $\abs c<\sqrt{2\eps}$, in polar coordinates, it integrates to $\frac14\sqrt{2\eps}\int_0^{\pi/2}(\cos\theta\sin\theta)^{-1/2}\dd\theta=\frac14\sqrt{2\eps}\,A_4$, step~3 again.
\end{computation}

\paragraph{What one had to notice.} Two things. Homogeneity makes every sublevel set a scaled copy of one set, so the volume is an exact power whose exponent is the number of coordinates divided by the degree: $\frac24$ here, $\frac22=\frac d2$ for the regular plane. And in its coefficients the quartic cup is the regular plane on a quadrant: the fold $a\mapsto a^2$ piles the prior up at the truth, and a parametrization that folds onto the truth pays less.\sidenote{The constant has a classical name: $\int_0^{\pi/2}(1-\frac12\sin^2\theta)^{-1/2}\dd\theta=\Gamma(\frac14)^2/(4\sqrt\pi)=1.8541$ is the complete elliptic integral of the first kind at modulus $1/\sqrt2$, and $C=\frac{\sqrt2}2\times1.8541$.} \Cref{fig:volume:shapes} shows the sets: the regular plane's disc shrinks like $\eps^{1/2}$ in each direction, the cup like $\eps^{1/4}$, the quartic model $K=\frac12(a^2+b^4)$ like $\eps^{1/2}$ in $a$ and $\eps^{1/4}$ in $b$, and the product model's cross keeps arms of full length that only thin.

\begin{figure}[t]
\begin{fullwidth}
\centering
\begin{tikzpicture}
\begin{groupplot}[group style={group size=4 by 1,horizontal sep=5mm,yticklabels at=edge left},
  splot,scale only axis,width=33mm,height=33mm,xmin=-1,xmax=1,ymin=-1,ymax=1,
  xtick={-1,0,1},ytick={-1,0,1},xlabel={$a$},title style={yshift=-1ex}]
\nextgroupplot[ylabel={$b$},title={regular plane}]
\addplot[fill=regular!14,draw=regular!70,line width=0.4pt] table[x=x1,y=y1] {figures/data/04-volume-shapes-reg.dat} -- cycle;
\addplot[fill=regular!34,draw=regular!80,line width=0.4pt] table[x=x2,y=y2] {figures/data/04-volume-shapes-reg.dat} -- cycle;
\addplot[fill=regular!65,draw=regular,line width=0.4pt] table[x=x3,y=y3] {figures/data/04-volume-shapes-reg.dat} -- cycle;
\nextgroupplot[title={quartic cup}]
\addplot[fill=singular!14,draw=singular!70,line width=0.4pt] table[x=x1,y=y1] {figures/data/04-volume-shapes-cup.dat} -- cycle;
\addplot[fill=singular!34,draw=singular!80,line width=0.4pt] table[x=x2,y=y2] {figures/data/04-volume-shapes-cup.dat} -- cycle;
\addplot[fill=singular!65,draw=singular,line width=0.4pt] table[x=x3,y=y3] {figures/data/04-volume-shapes-cup.dat} -- cycle;
\nextgroupplot[title={quartic model}]
\addplot[fill=singular!14,draw=singular!70,line width=0.4pt] table[x=x1,y=y1] {figures/data/04-volume-shapes-quart.dat} -- cycle;
\addplot[fill=singular!34,draw=singular!80,line width=0.4pt] table[x=x2,y=y2] {figures/data/04-volume-shapes-quart.dat} -- cycle;
\addplot[fill=singular!65,draw=singular,line width=0.4pt] table[x=x3,y=y3] {figures/data/04-volume-shapes-quart.dat} -- cycle;
\nextgroupplot[title={product model}]
\addplot[fill=singular!14,draw=singular!70,line width=0.4pt] table[x=x1,y=y1] {figures/data/04-volume-shapes-prod.dat} -- cycle;
\addplot[fill=singular!34,draw=singular!80,line width=0.4pt] table[x=x2,y=y2] {figures/data/04-volume-shapes-prod.dat} -- cycle;
\addplot[fill=singular!65,draw=singular,line width=0.4pt] table[x=x3,y=y3] {figures/data/04-volume-shapes-prod.dat} -- cycle;
\end{groupplot}
\end{tikzpicture}
\end{fullwidth}
\caption{The sets $\{K<\eps\}$ in the prior box $[-1,1]^2$ for $\eps=10^{-1}$, $10^{-2}$, $10^{-3}$ (light to dark), for $K=\frac12(a^2+b^2)$, $\frac12(a^4+b^4)$, $\frac12(a^2+b^4)$ and $\frac12a^2b^2$. Their prior masses are $0.157$, $0.0157$, $0.0016$; $0.415$, $0.131$, $0.041$; $0.261$, $0.046$, $0.008$; and $0.807$, $0.418$, $0.184$. Script: \texttt{figures/scripts/04-volume-shapes.py}.}\label{fig:volume:shapes}
\end{figure}

\paragraph{Every sum of powers.} Scaling works for every sum of powers of separate coordinates. Let $K=\frac12\sum_{j=1}^dw_j^{2k_j}$ with positive integers $k_j$, under the uniform prior on $[-1,1]^d$, and let $F(u)=\sum_ju_j^{2k_j}$ and $\lambda=\sum_j\frac1{2k_j}$. Scaling gives $\vol\{F<t\}=t^\lambda\vol\{F<1\}$, and computing $\int_{\R^d}e^{-F}$ both as a product of one-dimensional integrals and by slicing along the values of $F$ gives $\vol\{F<1\}\,\Gamma(\lambda+1)=\prod_j2\Gamma(1+\frac1{2k_j})$. Since $\{K<\eps\}=\{F<2\eps\}$ lies in the box for $\eps\le\frac12$,
\begin{equation}\label{eq:volume:sums}
V(\eps)=\frac{\prod_{j=1}^d\Gamma\bigl(1+\frac1{2k_j}\bigr)}{\Gamma(1+\lambda)}\,(2\eps)^{\lambda},\qquad\lambda=\sum_{j=1}^d\frac1{2k_j}\qquad\Bigl(0<\eps\le\frac12\Bigr).
\end{equation}
The exponents add: each coordinate contributes the reciprocal of its degree. The quartic model has $\lambda=\frac12+\frac14=\frac34$, the sextic model $\frac12(a^2+b^6)$ of \cref{ch:freeenergy} has $\frac23$, and the mixed cup $\frac12(a^4+b^6)$ of \cref{ch:newton} has $\frac5{12}$. Products need a second trick, the one with which \cref{ch:coefficient} found the density of a product of uniform variables.

\begin{computation}{Products in logarithmic coordinates}\label{comp:volume:log}
\step \emph{Logarithmic coordinates turn a uniform prior into exponentials.} If $a$ is uniform on $[-1,1]$, then $T_a=\log\frac1{\abs a}$ has $\Prob(T_a>t)=\Prob(\abs a<e^{-t})=e^{-t}$ for $t\ge0$: it is exponential with mean $1$. Under the uniform prior on $[-1,1]^2$, $T_a$ and $T_b$ are independent.
\step \emph{A product's sublevel set becomes a half-plane.} For the product model, with $s=\sqrt{2\eps}$ and $\ell=\log\frac1s$, $\{\frac12a^2b^2<\eps\}=\{T_a+T_b>\ell\}$. The sum of two independent standard exponentials has the tail $e^{-\ell}(1+\ell)$, so for $\eps\le\frac12$
\[
V(\eps)=s\Bigl(1+\log\frac1s\Bigr)\sim\frac1{\sqrt2}\,\eps^{1/2}\log\frac1\eps\qquad(\eps\to0):
\]
$\lambda=\frac12$ and $m=2$. The logarithm comes from the tie between two exponentials with the same rate.
\step \emph{Unequal rates: the slower one wins.} For the cusp model $y=ab^2x$, $K=\frac12a^2b^4$ and $\{K<\eps\}=\{T_a+2T_b>\ell\}$, where $2T_b$ is exponential with mean $2$. Conditioning on $T_a$,
\[
\Prob(T_a+2T_b>\ell)=e^{-\ell}+\int_0^\ell e^{-t}e^{-(\ell-t)/2}\dd t=2e^{-\ell/2}-e^{-\ell},
\]
so $V(\eps)=2(2\eps)^{1/4}-(2\eps)^{1/2}$: $\lambda=\frac14$, $m=1$, and the faster exponential leaves a correction of relative size $\eps^{1/4}$.
\step \emph{The rule.} With $T_j=\log\frac1{\abs{w_j}}$, a prior factor $\abs{w_j}^{h_j}$ makes $T_j$ exponential with rate $h_j+1$, and $\prod_jw_j^{2k_j}<\eps$ when $\sum_j2k_jT_j>\log\frac1\eps$, a sum of exponentials with rates $\frac{h_j+1}{2k_j}$. Its tail decays at the smallest rate, and each tie adds a power of $\log\frac1\eps$: this is \cref{prop:volume:monomial}.
\step \emph{State densities and zeta functions.} For the product model $v(t)=\log\frac1{2t}\big/(2\sqrt{2t})$ on $(0,\frac12)$, and $\zeta(z)=2^{-z}\bigl(\frac12\int_{-1}^1\abs a^{2z}\dd a\bigr)^2=2^{-z}(2z+1)^{-2}$: a double pole at $-\frac12$, with $(z+\frac12)^2\zeta(z)\to\frac{\sqrt2}4=0.3536$. For the cusp, $\zeta(z)=2^{-z}(2z+1)^{-1}(4z+1)^{-1}$: a simple pole at $-\frac14$, and one at $-\frac12$ that gives the correction of step~3.
\end{computation}

\begin{table}[t]
\centering\small
\begin{tabular}{@{}lllccS[table-format=1.4]c@{}}
\toprule
model & $K$ & $V(\eps)$ for $\eps\le\frac12$ & $\lambda$ & $m$ & {$C$} & $\frac d2$\\
\midrule
regular plane & $\frac12(a^2+b^2)$ & $\frac\pi2\,\eps$ & $1$ & $1$ & 1.5708 & $1$\\
quartic cup & $\frac12(a^4+b^4)$ & $1.3110\,\eps^{1/2}$ & $\frac12$ & $1$ & 1.3110 & $1$\\
quartic model & $\frac12(a^2+b^4)$ & $1.4699\,\eps^{3/4}$ & $\frac34$ & $1$ & 1.4699 & $1$\\
sextic model & $\frac12(a^2+b^6)$ & $1.4457\,\eps^{2/3}$ & $\frac23$ & $1$ & 1.4457 & $1$\\
mixed cup & $\frac12(a^4+b^6)$ & $1.2662\,\eps^{5/12}$ & $\frac5{12}$ & $1$ & 1.2662 & $1$\\
\addlinespace
product model & $\frac12a^2b^2$ & $s(1+\ell)$ & $\frac12$ & $2$ & 0.7071 & $1$\\
cusp model & $\frac12a^2b^4$ & $2s^{1/2}-s$ & $\frac14$ & $1$ & 2.3784 & $1$\\
cubic cusp & $\frac12a^2b^6$ & $\frac12(3s^{1/3}-s)$ & $\frac16$ & $1$ & 1.6837 & $1$\\
chain, $L=3$ & $\frac12a_1^2a_2^2a_3^2$ & $s(1+\ell+\frac12\ell^2)$ & $\frac12$ & $3$ & 0.1768 & $\frac32$\\
\bottomrule
\end{tabular}
\caption{Volumes of the running models under the uniform prior on $[-1,1]^d$, exact for $\eps\le\frac12$, with $s=\sqrt{2\eps}$ and $\ell=\log\frac1s$; $C$ is the constant in $V(\eps)\sim C\eps^\lambda(\log\frac1\eps)^{m-1}$. The first five are sums of powers, \eqref{eq:volume:sums}; the last four are monomials (\cref{comp:volume:log}, and \cref{prop:volume:monomial} for $C$). Every constant is checked by quadrature in \texttt{figures/scripts/04-volume-checks.py}.}\label{tab:volume:running}
\end{table}

\Cref{tab:volume:running} collects the two rules that the rest of Part~II elaborates: for a sum of powers of separate coordinates the exponents add; for a product the smallest exponent wins, and ties add powers of the logarithm. The quartic cup and the product model have the same learning coefficient $\frac12$ for different reasons: the cup because each coordinate contributes $\frac14$, the product because a small product needs only one small factor. Every analytic loss can be brought, piece by piece, to a monomial times a positive function by a change of coordinates (Hironaka's theorem, \cref{thm:toolkit:hironaka}); \cref{ch:blowups} does this by hand.

\section{The theorem: from the volume to the free energy}\label{sec:volume:abelian}

\begin{definition}[Volume, state density, zeta function, partition function]\label[definition]{def:volume:volume}
Let $W\subseteq\R^d$ be measurable, $\varphi$ a probability density on $W$, and $K\colon W\to[0,\infty)$ measurable. The volume function is $V(\eps)=\varphi\{w:K(w)<\eps\}$ for $\eps>0$. If $\varphi\{K=0\}=0$ and the law of $K(w)$ under $w\sim\varphi$ has a density $v$ on $(0,\infty)$, then $v$ is the \term{state density}\index{state density}, and $V(\eps)=\int_0^\eps v(t)\dd t$. The \term{zeta function}\index{zeta function} is
\[
\zeta(z)=\int_{\{K>0\}}K(w)^z\varphi(w)\dd w
\]
at the complex $z$ for which the integral converges absolutely. The \term{partition function}\index{partition function} is $Z(s)=\int_We^{-sK(w)}\varphi(w)\dd w$ and the \term{population free energy}\index{free energy!population} is $\bar F(s)=-\log Z(s)$, for $s>0$.
\end{definition}

All four depend only on the law of $K$ under the prior: $V$ is its distribution function, $\zeta$ its Mellin transform (\cref{def:toolkit:mellin}), and $Z$ its Laplace transform.

\begin{lemma}[Slicing by the value of the loss]\label[lemma]{lem:volume:layercake}
In the setting of \cref{def:volume:volume}, let $g\colon(0,\infty)\to[0,\infty)$ be measurable and $G(t)=\int_t^\infty g(u)\dd u$ for $t\ge0$. Then
\[
\int_WG\bigl(K(w)\bigr)\varphi(w)\dd w=\int_0^\infty g(t)\,V(t)\dd t .
\]
Both sides may be infinite. In particular $Z(s)=s\int_0^\infty e^{-st}V(t)\dd t$ for $s>0$, and $\int_WK^{-\sigma}\varphi=\sigma\int_0^\infty t^{-\sigma-1}V(t)\dd t$ for $\sigma>0$, with $K^{-\sigma}=\infty$ where $K=0$.
\end{lemma}

\begin{proof}
For each $w$, $G(K(w))=\int_0^\infty g(t)\ind\{K(w)<t\}\dd t$. Integrate against $\varphi$ and exchange the integrals, which Tonelli's theorem allows because the integrand is nonnegative; the inner integral $\int_W\ind\{K(w)<t\}\varphi(w)\dd w$ is $V(t)$. For the special cases take $g(t)=se^{-st}$, with $G(t)=e^{-st}$, and $g(t)=\sigma t^{-\sigma-1}$, with $G(t)=t^{-\sigma}$ for $t>0$ and $G(0)=\infty$.
\end{proof}

\begin{lemma}[From the volume to the free energy]\label[lemma]{lem:volume:abelian}
In the setting of \cref{def:volume:volume}, suppose that
\[
V(\eps)\sim C\eps^\lambda\Bigl(\log\frac1\eps\Bigr)^{m-1}\qquad(\eps\downarrow0)
\]
for constants $C>0$ and $\lambda>0$ and an integer $m\ge1$. Then, as $s\to\infty$ through the real numbers,
\begin{equation}\label{eq:volume:abelian}
\int_We^{-sK(w)}\varphi(w)\dd w\sim C\,\Gamma(\lambda+1)\,s^{-\lambda}(\log s)^{m-1},
\end{equation}
and therefore $\bar F(s)=\lambda\log s-(m-1)\log\log s-\log\bigl(C\,\Gamma(\lambda+1)\bigr)+o(1)$.
\end{lemma}

\begin{proof}
The plan is to write the partition function as an integral of $V$ against an exponential kernel, to rescale so that $V$ is evaluated where the hypothesis applies, and to pass to the limit by dominated convergence. By \cref{lem:volume:layercake} and the substitution $t=x/s$,
\[
\frac{Z(s)}{s^{-\lambda}(\log s)^{m-1}}=\int_0^\infty e^{-x}R_s(x)\dd x,\qquad R_s(x)=\frac{V(x/s)}{s^{-\lambda}(\log s)^{m-1}}.
\]
\emph{Pointwise limit.} Fix $x>0$. As $s\to\infty$, $x/s\to0$ and $\log\frac sx/\log s\to1$, so the hypothesis gives $R_s(x)\to Cx^\lambda$.

\emph{Domination.} Choose $t_0\in(0,1)$ such that $V(t)\le2Ct^\lambda(\log\frac1t)^{m-1}$ for $0<t<t_0$, and let $s\ge e$. If $x<st_0$, then $R_s(x)\le2Cx^\lambda\bigl(\log\frac sx/\log s\bigr)^{m-1}$, and $\log\frac sx/\log s=1+\log\frac1x/\log s$ is at most $1$ when $x\ge1$ and at most $1+\log\frac1x$ when $x<1$, because $\log s\ge1$. If $x\ge st_0$, then $V\le1$ and $(\log s)^{m-1}\ge1$ give $R_s(x)\le s^\lambda\le(x/t_0)^\lambda$. So for every $x>0$
\[
R_s(x)\le2Cx^\lambda\Bigl(1+\log^+\frac1x\Bigr)^{m-1}+\Bigl(\frac x{t_0}\Bigr)^\lambda,
\]
with $\log^+u=\max(\log u,0)$, and $e^{-x}$ times the right side is integrable on $(0,\infty)$ because $\lambda>0$. Dominated convergence gives $\int_0^\infty e^{-x}R_s(x)\dd x\to C\int_0^\infty x^\lambda e^{-x}\dd x=C\,\Gamma(\lambda+1)$, which is \eqref{eq:volume:abelian}; taking logarithms gives the statement about $\bar F$.
\end{proof}

The lemma is step~5 of \cref{comp:volume:cup} in general, where $\sqrt2\,\Gamma(\frac54)^2=1.1619=C\,\Gamma(\frac32)$. \Cref{ch:freeenergy} applies it, at real $s$, to the multiples $(1\pm\eta)^2K$ that sandwich the random loss (\cref{thm:freeenergy:linear}). Its converse is Karamata's Tauberian theorem (\cref{thm:toolkit:karamata}), so $V$ and $Z$ determine each other's leading terms.\sidenote{Theorems that pass from a function to its transform are called Abelian, and those that pass back, which need an extra condition such as monotonicity, Tauberian; \textcite{BinghamEtAl1987} is the standard reference.} When $m\ge2$ the limit is approached slowly: for the product model $Z(n)=\frac12\sqrt{\pi/(2n)}\,(\log n+\gamma+\log2)$ up to an exponentially small error, with $\gamma$ Euler's constant, so its ratio to the lemma's prediction is $1.069$ at $n=10^8$ and reaches $1.01$ only at $n\approx10^{55}$.

The next result computes the volume of every monomial times a positive factor, under a prior of the same kind; by Hironaka's theorem it is the local building block of every learning coefficient in the book.

\begin{proposition}[Volumes of normal-crossing functions]\label[proposition]{prop:volume:monomial}
Let $B=\prod_{j=1}^d[-r_j,r_j]$ with $r_j>0$. Let $k_1,\dots,k_d$ and $h_1,\dots,h_d$ be nonnegative integers, not all $k_j$ zero, and let $a$ and $b$ be continuous functions on $B$ with $a>0$ and $b\ge0$. On $B$ put
\[
K(w)=a(w)\prod_{j=1}^dw_j^{2k_j},\qquad\varphi(w)=b(w)\prod_{j=1}^d\abs{w_j}^{h_j},
\]
and $V(\eps)=\int_{B\cap\{K<\eps\}}\varphi$. For $j$ with $k_j>0$ let $\lambda_j=\frac{h_j+1}{2k_j}$; let $\lambda=\min_j\lambda_j$, let $J$ be the set of indices with $\lambda_j=\lambda$, and let $m$ be its number of elements. Write $w'=(w_j)_{j\notin J}$, $B'=\prod_{j\notin J}[-r_j,r_j]$, and $(0_J,w')$ for the point whose coordinates in $J$ are $0$. Then $V(\eps)=(C+o(1))\,\eps^\lambda(\log\frac1\eps)^{m-1}$ as $\eps\downarrow0$, with
\begin{equation}\label{eq:volume:monomialC}
C=\frac{1}{(m-1)!\,\lambda\prod_{j\in J}k_j}\int_{B'}\frac{b(0_J,w')}{a(0_J,w')^\lambda}\prod_{j\notin J}\abs{w_j}^{h_j-2k_j\lambda}\dd w',
\end{equation}
a finite integral, read as $b(0)a(0)^{-\lambda}$ when $J$ contains every index. If $b(0)>0$, then $C>0$ and $V(\eps)\sim C\eps^\lambda(\log\frac1\eps)^{m-1}$.
\end{proposition}

\begin{proof}
The plan is to reduce to the positive orthant of the unit cube, to compute the volume exactly in the coordinates of $J$ by the logarithmic coordinates of \cref{comp:volume:log}, to integrate the other coordinates by dominated convergence with $a$ and $b$ frozen at $w_J=0$, and to control the error of freezing by continuity. Write $L=\log\frac1\eps$ and $\pi(u')=\prod_{j\notin J}u_j^{h_j-2k_j\lambda}$.

\emph{Step 1: reductions.} The coordinate hyperplanes are null sets, so $V$ is the sum of the volumes of the $2^d$ orthants. A reflection $w\mapsto(\sigma_jw_j)_j$ maps an orthant onto $\prod_j[0,r_j]$ and changes only $a$ and $b$, into $a(\sigma\,\cdot)$ and $b(\sigma\,\cdot)$. The substitution $w_j=r_ju_j$ maps $\prod_j[0,r_j]$ onto $[0,1]^d$ and keeps the form of $K$ and $\varphi\dd w$, with $a$ and $b$ multiplied by $\prod_jr_j^{2k_j}$ and $\prod_jr_j^{h_j+1}$; the formula for $C$ does not change, because $h_j+1=2k_j\lambda$ for $j\in J$ and $r_j^{h_j+1-2k_j\lambda}u_j^{h_j-2k_j\lambda}\dd u_j=w_j^{h_j-2k_j\lambda}\dd w_j$ for $j\notin J$. It therefore suffices to prove, for all continuous $a>0$ and $b\ge0$ on $[0,1]^d$, that $V(\eps)=(C_++o(1))\eps^\lambda L^{m-1}$ with
\begin{equation}\label{eq:volume:orthant}
C_+=c_0\int_{[0,1]^{J^c}}\frac{b(0_J,u')}{a(0_J,u')^\lambda}\,\pi(u')\dd u',\qquad c_0=\frac{1}{2^m(m-1)!\,\lambda\prod_{j\in J}k_j}.
\end{equation}
Summing over the orthants then gives \eqref{eq:volume:monomialC}: at $u_J=0$ the signs of the coordinates in $J$ do not change the integrand, which gives the factor $2^m$, and the other signs reassemble the integral over $B'$. The integral is finite because $h_j-2k_j\lambda>-1$ for $j\notin J$ (if $k_j>0$ this says $\lambda<\lambda_j$; if $k_j=0$ it says $h_j>-1$), $a$ is bounded below and $b$ above.

\emph{Step 2: the coordinates in $J$, exactly.} For $\eta>0$ let $G(\eta)=\int_{[0,1]^J}\ind\{\prod_{j\in J}u_j^{2k_j}<\eta\}\prod_{j\in J}u_j^{h_j}\dd u_J$. For $j\in J$ substitute $u_j=e^{-t_j/(h_j+1)}$, $t_j\ge0$: then $u_j^{h_j}\dd u_j$ becomes $(h_j+1)^{-1}e^{-t_j}\dd t_j$, and $u_j^{2k_j}=e^{-t_j/\lambda}$ because $\frac{2k_j}{h_j+1}=\frac1\lambda$. The condition reads $\sum_{j\in J}t_j>\lambda\log\frac1\eta$; under the densities $e^{-t_j}$ the sum has the Gamma density $t^{m-1}e^{-t}/(m-1)!$, whose tail at $x\ge0$ is $e^{-x}\sum_{i<m}x^i/i!$ (both equal $1$ at $x=0$ and have the derivative $-x^{m-1}e^{-x}/(m-1)!$). With $c_J=\prod_{j\in J}(h_j+1)^{-1}$,
\begin{equation}\label{eq:volume:gammatail}
G(\eta)=c_J\,\eta^\lambda\sum_{i=0}^{m-1}\frac{(\lambda\log\frac1\eta)^i}{i!}\quad(0<\eta\le1),\qquad G(\eta)=c_J\quad(\eta\ge1).
\end{equation}

\emph{Step 3: frozen coefficients.} Let $F\ge0$ be bounded and measurable and $\alpha$ measurable with values in $[\alpha_-,\alpha_+]\subset(0,\infty)$, both functions of $u'=(u_j)_{j\notin J}$ only, and let $M(u')=\prod_{j\notin J}u_j^{2k_j}$. By Fubini's theorem the volume of $\{\alpha M\prod_{j\in J}u_j^{2k_j}<\eps\}$ under $F(u')\prod_ju_j^{h_j}\dd u$ is
\[
I(\eps)=\int_{[0,1]^{J^c}}F(u')\prod_{j\notin J}u_j^{h_j}\;G\Bigl(\frac\eps{\alpha(u')M(u')}\Bigr)\dd u'
\]
(no integral when $J$ contains every index). We claim that $I(\eps)/(\eps^\lambda L^{m-1})\to c_0\int F\alpha^{-\lambda}\pi\dd u'$; note that $c_J\lambda^{m-1}/(m-1)!=c_0$. Fix $u'$ with $M(u')>0$, which excludes a null set, and let $\eta=\eps/(\alpha M)$. As $\eps\to0$, $\eta\to0$, and by \eqref{eq:volume:gammatail}, since $\log\frac1\eta=L+\log(\alpha M)$,
\[
\frac{G(\eta)}{\eps^\lambda L^{m-1}}=\frac{c_J}{(\alpha M)^\lambda}\sum_{i<m}\frac{\lambda^i}{i!}\frac{(\log\frac1\eta)^i}{L^{m-1}}\to\frac{c_0}{(\alpha M)^\lambda}.
\]
For domination let $L\ge\max(1,\log\alpha_+)$. If $\eta\le1$, then $0\le\log\frac1\eta\le2L$, because $M\le1$, so $G(\eta)\le c_Je^{2\lambda}\eta^\lambda L^{m-1}$; if $\eta>1$, then $G(\eta)=c_J\le c_J\eta^\lambda$. In both cases $G(\eta)/(\eps^\lambda L^{m-1})\le c_Je^{2\lambda}\alpha_-^{-\lambda}M(u')^{-\lambda}$, and $F\prod_{j\notin J}u_j^{h_j}M^{-\lambda}=F\pi$ is integrable by Step~1. Dominated convergence proves the claim, which for $F=\alpha=1$ is the proposition with $a=b=1$.

\emph{Step 4: unfreezing.} Write $u^\circ=(0_J,u')$, $a_-=\min a$ and $b_+=\max b$, and fix $\eta'\in(0,\frac12)$. By uniform continuity there is $\delta\in(0,1)$ such that $\abs{b(u)-b(u^\circ)}\le\eta'$ and $(1-\eta')a(u^\circ)\le a(u)\le(1+\eta')a(u^\circ)$ whenever $\max_{j\in J}u_j<\delta$. Split the cube into $N=\{\max_{j\in J}u_j<\delta\}$ and its complement.

The complement contributes $o(\eps^\lambda L^{m-1})$. If $u_{j_1}\ge\delta$ for some $j_1\in J$, then $K<\eps$ forces $\prod_{j\ne j_1}u_j^{2k_j}<\eps/(a_-\delta^{2k_{j_1}})$, so the contribution is at most $\sum_{j_1\in J}b_+V_{j_1}(\eps/(a_-\delta^{2k_{j_1}}))$, with $V_{j_1}$ the volume of $\prod_{j\ne j_1}u_j^{2k_j}$ on $[0,1]^{d-1}$ under $\prod_{j\ne j_1}u_j^{h_j}\dd u$. If no $k_j$ with $j\ne j_1$ is positive, $V_{j_1}(\eta)=0$ for $\eta\le1$. Otherwise Steps 2 and~3 with $F=\alpha=1$ apply to that monomial, whose smallest ratio is $\lambda$, attained $m-1$ times, when $m\ge2$, and some $\lambda'>\lambda$ when $m=1$; so the contribution is $O(\eps^\lambda L^{m-2})$ or $O(\eps^{\lambda'}L^{d})$, both $o(\eps^\lambda L^{m-1})$.

On $N$ we have $b\le b(u^\circ)+\eta'$ and $K\ge(1-\eta')a(u^\circ)\prod_ju_j^{2k_j}$, so the volume of $N\cap\{K<\eps\}$ is at most the frozen integral $I_+(\eps)$ with $F=b(u^\circ)+\eta'$ and $\alpha=(1-\eta')a(u^\circ)$. We also have $b\ge(b(u^\circ)-\eta')_+$ and $K\le(1+\eta')a(u^\circ)\prod_ju_j^{2k_j}$ on $N$, so that volume is at least the frozen integral $I_-(\eps)$ with $F=(b(u^\circ)-\eta')_+$ and $\alpha=(1+\eta')a(u^\circ)$, minus the part of $I_-(\eps)$ outside $N$, which is $o(\eps^\lambda L^{m-1})$ as above. Step~3 applied to $I_+$ and $I_-$ gives
\[
\limsup_{\eps\to0}\frac{V(\eps)}{\eps^\lambda L^{m-1}}\le\frac{c_0}{(1-\eta')^\lambda}\int\frac{b(u^\circ)+\eta'}{a(u^\circ)^\lambda}\,\pi\dd u',
\]
\[
\liminf_{\eps\to0}\frac{V(\eps)}{\eps^\lambda L^{m-1}}\ge\frac{c_0}{(1+\eta')^\lambda}\int\frac{(b(u^\circ)-\eta')_+}{a(u^\circ)^\lambda}\,\pi\dd u'.
\]
As $\eta'\downarrow0$ both bounds tend to $C_+$ by dominated convergence. Finally, if $b(0)>0$, then $b(0_J,w')>0$ for $w'$ near $0$, and the integral in \eqref{eq:volume:monomialC} is positive.
\end{proof}

The proof works verbatim for real $k_j\ge0$ and $h_j>-1$. The constants of \cref{tab:volume:running} are the proposition with $a=\frac12$ and $b=2^{-d}$. A check with non-constant factors, $K=(1+\frac12u_1+u_2^2)\,u_1^2u_2^4$ and $\varphi=(2+\sin u_1+u_2)\abs{u_1}$ on $[-1,1]^2$, gives $V(\eps)/\eps^{1/4}=5.29322$ at $\eps=10^{-8}$ and $5.29364$ at $10^{-16}$ by quadrature, against $C=5.29364$. On $[0,1]^d$ with $a=b=1$, $\lambda C_+$ is the total mass of the measure $\dd u^*$ of \textcite[Lemma~4]{Watanabe2013}. The volume also gives rules of \cref{ch:blowups} (\cref{prop:blowups:rules}) in one line: if $c_1K_1\le K_2\le c_2K_1$, then $V_1(\eps/c_2)\le V_2(\eps)\le V_1(\eps/c_1)$, so comparable losses share $\lambda$ and $m$ (\cref{ex:volume:rules}).

\section{The theorem: the volume and the zeta function}\label{sec:volume:zeta}

The book defines the learning coefficient through the zeta function. The volume function and the zeta function are two transforms of one law, so they carry the same information; this section says how.

\begin{definition}[Learning coefficient and multiplicity]\label[definition]{def:volume:lambda}
In the setting of \cref{def:volume:volume}, let $K$ be bounded and $\varphi\{K=0\}=0$. Let $\lambda>0$, and suppose that the integral defining $\zeta$ converges for $\Re z>-\lambda$ and that $\zeta$ extends to a meromorphic function on an open set containing the closed half-plane $\{\Re z\ge-\lambda\}$, whose only pole there is $-\lambda$, of order $m$. Then $\lambda$ is the \term{learning coefficient}\index{learning coefficient} of $K$ with respect to $\varphi$, and $m$ is its \term{multiplicity}\index{multiplicity}; $\lambda$ is also called the \term{real log canonical threshold}\index{real log canonical threshold}.
\end{definition}

For analytic losses (part~(iv) below) and every running model, $-\lambda$ is simply the largest pole. A pole of order $m$ means that near $-\lambda$, $\zeta(z)=(z+\lambda)^{-m}P(z)+h(z)$ with a polynomial $P$, a convergent power series $h$, and $P(-\lambda)\ne0$, the \emph{leading coefficient}.

\begin{theorem}[The volume and the zeta function]\label{thm:volume:zeta}
In the setting of \cref{def:volume:volume}, let $K$ be bounded, $\varphi\{K=0\}=0$, and $\lambda>0$.
\begin{enumerate}
\item \emph{(Convergence.)} The real $x$ with $\int_{\{K>0\}}K^x\varphi<\infty$ form an interval containing $[0,\infty)$; let $-\sigma_c\le0$ be its left end point. On $\Re z>-\sigma_c$ the integral defining $\zeta$ converges absolutely and $\zeta$ is holomorphic, and for $-\sigma_c<\Re z<0$
\begin{equation}\label{eq:volume:mellin}
\zeta(z)=-z\int_0^\infty t^{z-1}\,V(t)\dd t .
\end{equation}
\item \emph{(From the volume to the zeta function.)} If $V(\eps)\sim C\eps^\lambda(\log\frac1\eps)^{m-1}$ with $C>0$ and an integer $m\ge1$, then $\sigma_c=\lambda$, the integral diverges at $z=-\lambda$, and
\begin{equation}\label{eq:volume:zetalimit}
(z+\lambda)^m\zeta(z)\to\lambda\,C\,(m-1)!\qquad(z\downarrow-\lambda\text{ through real }z).
\end{equation}
Consequently, if $\zeta$ extends meromorphically to a neighbourhood of $-\lambda$, then $-\lambda$ is a pole of order exactly $m$ with leading coefficient $\lambda C(m-1)!$, and no pole lies in $\Re z>-\lambda$.
\item \emph{(From the zeta function to the volume.)} Conversely, if $\zeta$ satisfies the hypothesis of \cref{def:volume:lambda}, with leading coefficient $A$ at $-\lambda$, then $A>0$ and
\[
V(\eps)\sim\frac{A}{\lambda\,(m-1)!}\,\eps^\lambda\Bigl(\log\frac1\eps\Bigr)^{m-1}\qquad(\eps\downarrow0).
\]
\item \emph{(Analytic losses; statement.)} Let $W=\{w:\pi_1(w)\ge0,\dots,\pi_k(w)\ge0\}$ be compact with nonempty interior, with analytic $\pi_j$; let $\varphi=\varphi_1\varphi_2$ with $\varphi_1\ge0$ analytic and $\varphi_2>0$ infinitely differentiable; and let $K\ge0$ be analytic on a neighbourhood of $W$ with $\varphi\{K=0\}=0$. Then $\zeta$ extends to a meromorphic function on $\C$ whose poles are negative rational numbers. If it has a pole, then its largest pole $-\lambda$ and that pole's order $m$ satisfy the hypothesis of \cref{def:volume:lambda}, and $V(\eps)\sim C\eps^\lambda(\log\frac1\eps)^{m-1}$ with $C>0$.
\end{enumerate}
\end{theorem}

So the learning coefficient is the exponent of the volume, the multiplicity is one more than the power of its logarithm, and the leading coefficient is $\lambda C(m-1)!$: $0.6555$ for the quartic cup and $0.3536$ for the product model, as \cref{comp:volume:cup,comp:volume:log} found. Parts (i) and~(ii) are elementary; part~(iii) rests on a Tauberian theorem of \textcite{Delange1954}. Part~(iv) is stated: the continuation is the theorem of \textcite{BernsteinGelfand1969} and \textcite{Atiyah1970}, which Watanabe states for the zeta functions of statistical models \parencite[p.~3584]{Watanabe2010}, and given a resolution of singularities \cref{ch:blowups} derives the rest from \cref{prop:volume:monomial} (\cref{thm:blowups:pullback}).

\begin{proof}[Proof of (i) and (ii)]
The plan is to read convergence off \cref{lem:volume:layercake}, to obtain holomorphy from power series, and to prove \eqref{eq:volume:zetalimit} by an Abelian argument in logarithmic coordinates. Let $T=\max(1,\sup K)$.

\emph{(i).} If $x<y$ are real and $\int K^x\varphi<\infty$, then $K^y\le K^x+T^{\abs y}$ on $\{K>0\}$, because $K^y\le K^x$ where $K\le1$ and $K^y\le T^{\abs y}$ where $K>1$; so these $x$ form an interval, which contains $[0,\infty)$ because $K^x\le T^x$ for $x\ge0$. Let $\Re z_0>-\sigma_c$, and choose $r>0$ with $\Re z_0-r>-\sigma_c$. For $\abs{z-z_0}<r$ and $K(w)>0$,
\[
K^z=K^{z_0}e^{(z-z_0)\log K}=\sum_{k\ge0}\frac{(z-z_0)^k}{k!}\,K^{z_0}(\log K)^k,
\]
and the series of absolute values is $K^{\Re z_0}e^{\abs{z-z_0}\abs{\log K}}\le K^{\Re z_0-r}+K^{\Re z_0+r}$, which is integrable. So the series may be integrated term by term, and near $z_0$ the function $\zeta$ is a convergent power series in $z-z_0$: it is holomorphic. For \eqref{eq:volume:mellin} let $\Re z<0$. For $K(w)>0$, $K(w)^z=-z\int_{K(w)}^\infty t^{z-1}\dd t$; the integral of $\abs z\,t^{\Re z-1}\ind\{t>K(w)\}\varphi(w)$ over $t$ and $w$ is $\frac{\abs z}{\abs{\Re z}}\int K^{\Re z}\varphi<\infty$, so Fubini's theorem allows the exchange, and $\varphi\{0<K<t\}=V(t)$ because $\varphi\{K=0\}=0$.

\emph{(ii), convergence.} For $\sigma>0$, \cref{lem:volume:layercake} gives $\int K^{-\sigma}\varphi=\sigma\int_0^\infty t^{-\sigma-1}V(t)\dd t$, and the part over $[1,\infty)$ is at most $1$. Near $t=0$, $V(t)$ lies between $\frac C2t^\lambda(\log\frac1t)^{m-1}$ and $2Ct^\lambda(\log\frac1t)^{m-1}$, so the part over $(0,1]$ is finite exactly when $\sigma<\lambda$. Hence $\sigma_c=\lambda$, and the integral diverges at $-\lambda$.

\emph{(ii), the limit.} Let $0<\sigma<\lambda$ and $\delta=\lambda-\sigma$. The substitution $t=e^{-u}$ gives $\zeta(-\sigma)=\sigma\int_0^\infty e^{-\delta u}f(u)\dd u+\theta_\sigma$ with $f(u)=e^{\lambda u}V(e^{-u})$ and $0\le\theta_\sigma\le1$. By hypothesis $f(u)\sim Cu^{m-1}$ as $u\to\infty$, and $0\le f(u)\le e^{\lambda u}$. Given $\eta'>0$, choose $U\ge1$ with $\abs{f(u)-Cu^{m-1}}\le\eta'u^{m-1}$ for $u\ge U$. Since $\int_0^\infty e^{-\delta u}u^{m-1}\dd u=\Gamma(m)\delta^{-m}$,
\[
\Bigl\lvert\delta^m\int_U^\infty e^{-\delta u}f(u)\dd u-C\int_{\delta U}^\infty e^{-y}y^{m-1}\dd y\Bigr\rvert\le\eta'\,\Gamma(m),
\]
while the integral over $[0,U]$, multiplied by $\delta^m$, lies between $0$ and $\delta^mUe^{\lambda U}$. As $\delta\downarrow0$ the integral over $y\ge\delta U$ tends to $\Gamma(m)$ and $\delta^mUe^{\lambda U}$ to $0$; since $\eta'$ was arbitrary, $\delta^m\int_0^\infty e^{-\delta u}f(u)\dd u\to C\,\Gamma(m)$. Multiplying by $\sigma\to\lambda$, and noting $\delta^m\theta_\sigma\to0$, gives \eqref{eq:volume:zetalimit}.

\emph{(ii), the consequence.} If $\zeta$ extends meromorphically near $-\lambda$, its expansion there is $\sum_{k\ge-p}c_k(z+\lambda)^k$ with $c_{-p}\ne0$, and $(z+\lambda)^m\zeta(z)$ behaves like $c_{-p}(z+\lambda)^{m-p}$, which has a finite nonzero limit only if $p=m$. By~(i) no pole lies in $\Re z>-\lambda$.
\end{proof}

\begin{proof}[Proof of (iii)]
The plan is to apply Delange's Tauberian theorem to a monotone function built from the law of $\log\frac1K$, and to return to $V$ by an integration by parts. Replacing $K$ by $K/T$, $T=\max(1,\sup K)$, multiplies $\zeta(z)$ by $T^{-z}$ and $A$ by $T^\lambda$, and replaces $V(\eps)$ by $V(T\eps)$; the claim is unchanged, so we may assume $K\le1$. Let $\mu$ be the law of $U=\log\frac1{K(w)}$ under $w\sim\varphi$, a law on $[0,\infty)$; then $\zeta(z)=\int e^{-zu}\mu(\dd u)$ for $\Re z>-\lambda$, and $V(e^{-t})=\mu((t,\infty))$.

\emph{A monotone function and its transform.} Fix $\theta>0$ and, for $t\ge0$, let
\[
\alpha(t)=\int_{[0,t]}e^{(\lambda+\theta)u}\mu(\dd u),
\]
a nondecreasing nonnegative function. For real $s>\theta$, Tonelli's theorem gives
\[
\int_0^\infty e^{-st}\alpha(t)\dd t=\int e^{(\lambda+\theta)u}\Bigl(\int_u^\infty e^{-st}\dd t\Bigr)\mu(\dd u)=\frac{\zeta(s-\lambda-\theta)}s,
\]
which is finite because $s-\lambda-\theta>-\lambda$. So the Laplace transform $f$ of $\alpha$ converges for $\Re s>\theta$ and equals $\zeta(s-\lambda-\theta)/s$ there. By hypothesis $f$ extends holomorphically to every point of the line $\Re s=\theta$ other than $\theta$, and near $\theta$, $f(s)=(s-\theta)^{-m}g(s)+h(s)$ with $g(s)=P(s-\lambda-\theta)/s$, $h$ holomorphic at $\theta$, and $g(\theta)=A/\theta>0$, because $A\ne0$ and $(z+\lambda)^m\zeta(z)>0$ for real $z>-\lambda$. These are the hypotheses of \textcite[Th\'eor\`eme~III]{Delange1954}, with $a=\theta>0$ and $\omega=m$, not an integer $\le0$; its conclusion is $\alpha(t)\sim Be^{\theta t}t^{m-1}$ as $t\to\infty$, with $B=\frac{A}{\theta\,(m-1)!}$.

\emph{Back to the volume.} Since $e^{-(\lambda+\theta)u}=(\lambda+\theta)\int_u^\infty e^{-(\lambda+\theta)y}\dd y$, Tonelli's theorem and $\int_{(t,y]}\dd\alpha=\alpha(y)-\alpha(t)$ give
\[
V(e^{-t})=\int_{(t,\infty)}e^{-(\lambda+\theta)u}\dd\alpha(u)=(\lambda+\theta)\int_t^\infty e^{-(\lambda+\theta)y}\alpha(y)\dd y-e^{-(\lambda+\theta)t}\alpha(t).
\]
As $t\to\infty$, the first term is $\frac{\lambda+\theta}\lambda Be^{-\lambda t}t^{m-1}(1+o(1))$, because $\int_t^\infty e^{-\lambda y}y^{m-1}\dd y\sim e^{-\lambda t}t^{m-1}/\lambda$, and the second is $Be^{-\lambda t}t^{m-1}(1+o(1))$. Their difference is $B(\frac\theta\lambda+o(1))e^{-\lambda t}t^{m-1}$, which with $\eps=e^{-t}$ is the claim.
\end{proof}

The proof never bounded $\zeta$ along vertical lines; positivity of the law of $K$ did that work, through Delange's theorem.\sidenote{The hypothesis that the integral converges for $\Re z>-\lambda$ can be dropped. After scaling to $K\le1$: if $\zeta$ were holomorphic at the left end point $-\sigma_c>-\lambda$ of its half-plane of convergence, its Taylor series at a nearby real point $x$, with coefficients $\int K^x(\log K)^k\varphi/k!$, would converge a little to the left of $-\sigma_c$ with terms of one sign there, and Tonelli's theorem would make the integral converge there too.} The more familiar route, Mellin inversion, shows where the logarithms come from but needs such bounds.

\begin{heuristic}[Why a pole of order $m$ gives the $(m-1)$st power of the logarithm]
By \eqref{eq:volume:mellin}, $-\zeta(z)/z$ is the Mellin transform of $V$, so $V(t)=\frac1{2\pi i}\int_{c-i\infty}^{c+i\infty}t^{-z}\,\frac{-\zeta(z)}z\dd z$ for $-\lambda<c<0$. Move the line of integration left, past $-\lambda$. Since $t^{-z}=t^\lambda\sum_k(\log\frac1t)^k(z+\lambda)^k/k!$, the residue at a pole $A(z+\lambda)^{-m}+\dots$ is $t^\lambda$ times a polynomial of degree $m-1$ in $\log\frac1t$ with leading term $\frac A\lambda(\log\frac1t)^{m-1}/(m-1)!$, and the new line contributes $O(t^{\lambda+\delta})$. \Textcite[\S\S4.2--4.4]{Watanabe2009} proceeds this way for normal-crossing functions, from the poles of the zeta function to the expansion of the state density.
\end{heuristic}

\paragraph{Where the zeta function comes from.} In his address to the 1954 International Congress of Mathematicians in Amsterdam, I.\,M.\ Gel'fand asked whether, for a polynomial $P$ and a test function $\phi$, the integral $\int_{\{P>0\}}P^\lambda\phi\dd x$ is a meromorphic function of $\lambda$ with poles in finitely many arithmetic progressions, and asked for the residues \parencite[\S3, pp.~262--263]{Gelfand1957}. It was answered with Hironaka's resolution of singularities \parencite{Hironaka1964a} by I.\,N.\ Bernstein and S.\,I.\ Gel'fand \parencite{BernsteinGelfand1969}, and independently by \textcite{Atiyah1970}, who used the result to divide distributions by analytic functions: after a resolution the integral is a finite sum of integrals of monomials times smooth functions, whose poles are negative rationals. \Textcite{Bernstein1972} proved the continuation for polynomials without resolution, from an identity $b(z)P^z=D(z)P^{z+1}$ with a polynomial $b$ and a differential operator $D$ polynomial in $z$; $b$ is now called the \term{Bernstein--Sato polynomial}\index{Bernstein--Sato polynomial}, and \textcite{Kashiwara1976} proved that its roots are negative rationals. \Textcite{Watanabe1999,Watanabe2001a} brought the zeta function into statistics.

\section{Measuring it}\label{sec:volume:measuring}

The script \texttt{04-volume-mc.py} estimates the volumes of the regular plane and of the quartic cup, product, quartic and sextic models, under the uniform prior on $[-1,1]^2$, by plain Monte Carlo: it draws $N$ points from the prior with the kit's \texttt{volume\_fraction}, counts those with $K<\eps$ at the fifteen levels $\eps=10^{-1},10^{-1.5},\dots,10^{-8}$, and reports each fraction with its exact Clopper--Pearson $95\%$ interval (\cref{sec:toolkit:intervals}). The count is binomial with mean $NV(\eps)$, so the relative standard error is about $(NV)^{-1/2}$ and $10\%$ accuracy needs $N\approx100/V(\eps)$: at $\eps=10^{-8}$, $7.2\times10^4$ draws for the product model, $7.6\times10^5$ for the cup, $1.5\times10^7$ for the sextic model, $6.8\times10^7$ for the quartic model and $6.4\times10^9$ for the regular plane. The script uses $10^7$, $10^7$, $3\times10^7$, $10^8$ and $10^8$ draws (seeds $401$--$405$), at least $153$ inside every level for the singular models; it and \texttt{04-volume-zeta.py} each run in under a minute.

\begin{figure}[t]
\begin{fullwidth}
\centering
\begin{tikzpicture}
\begin{groupplot}[group style={group size=2 by 1,horizontal sep=17mm},splot,width=80mm,height=64mm]
\nextgroupplot[xmode=log,ymode=log,xmin=4e-9,xmax=0.25,ymin=1e-12,ymax=2,xtick={1e-8,1e-6,1e-4,1e-2},
  ytick={1e-12,1e-9,1e-6,1e-3,1},xlabel={level $\eps$},ylabel={volume $V(\eps)$},
  title={(a) the volume function},legend pos=south east,legend style={font=\scriptsize}]
\addplot[regularpred] table[x=eps,y=reg] {figures/data/04-volume-exact.dat};
\addplot[theory] table[x=eps,y=cup] {figures/data/04-volume-exact.dat};
\addplot[measured3,ebars] table[x=eps,y=V,y error plus=ep,y error minus=em] {figures/data/04-volume-mc-reg.dat};
\addplot[measured,ebars] table[x=eps,y=V,y error plus=ep,y error minus=em] {figures/data/04-volume-mc-cup.dat};
\addplot[measured2,ebars] table[x=eps,y=V,y error plus=ep,y error minus=em] {figures/data/04-volume-mc-prod.dat};
\addplot[only marks,mark=diamond*,mark size=2pt,color=singulardeep,mark options={fill=white,line width=0.6pt},ebars] table[x=eps,y=V,y error plus=ep,y error minus=em] {figures/data/04-volume-mc-quart.dat};
\addplot[only marks,mark=triangle*,mark size=2pt,color=muted,mark options={fill=white,line width=0.6pt},ebars] table[x=eps,y=V,y error plus=ep,y error minus=em] {figures/data/04-volume-mc-sext.dat};
\addplot[theory,forget plot] table[x=eps,y=prod] {figures/data/04-volume-exact.dat};
\addplot[theory,forget plot] table[x=eps,y=quart] {figures/data/04-volume-exact.dat};
\addplot[theory,forget plot] table[x=eps,y=sext] {figures/data/04-volume-exact.dat};
\legend{regular prediction $\frac\pi2\eps$,singular: exact volumes,regular plane,quartic cup,product model,quartic model,sextic model}
\nextgroupplot[xmin=-0.5,xmax=0.0,ymin=0,ymax=0.72,xtick={-0.5,-0.4,-0.3,-0.2,-0.1,0},
  ytick={0,0.1,0.2,0.3,0.4,0.5,0.6,0.7},xlabel={$z$},ylabel={$(z+\frac12)^m\,\zeta(z)$},
  title={(b) the zeta function near its pole},legend pos=south east,legend style={font=\scriptsize}]
\addplot[guide,forget plot] coordinates {(-0.25,0) (-0.25,0.72)};
\addplot[theory] table[x=z,y=cupscaled] {figures/data/04-volume-zeta-exact.dat};
\addplot[theory2] table[x=z,y=prodscaled] {figures/data/04-volume-zeta-exact.dat};
\addplot[measured,ebars] table[x=z,y=cup,y error plus=cuphalf,y error minus=cuphalf] {figures/data/04-volume-zeta-mc.dat};
\addplot[measured2,ebars] table[x=z,y=prod,y error plus=prodhalf,y error minus=prodhalf] {figures/data/04-volume-zeta-mc.dat};
\addplot[only marks,mark=x,mark size=2.5pt,color=singulardeep,line width=0.8pt] coordinates {(-0.5,0.655514) (-0.5,0.353553)};
\legend{cup ($m=1$): exact,product ($m=2$): exact,cup: Monte Carlo,product: Monte Carlo,limits $\lambda C(m-1)!$}
\end{groupplot}
\end{tikzpicture}
\end{fullwidth}
\caption{Monte Carlo against exact values under the uniform prior on $[-1,1]^2$. (a) The fraction of $N$ prior draws with $K<\eps$ at fifteen levels, with exact Clopper--Pearson $95\%$ intervals (mostly smaller than the marks), for the regular plane ($N=10^8$, seed $401$), the quartic cup ($10^7$, $402$), the product model ($10^7$, $403$), the quartic model ($10^8$, $404$) and the sextic model ($3\times10^7$, $405$). Solid bronze: exact volumes (\cref{tab:volume:running}). Dashed slate: the regular prediction, the $\eps^{d/2}$ law, which the regular plane follows exactly. (b) $(z+\frac12)^m\zeta(z)$ for the quartic cup ($m=1$) and the product model ($m=2$). Bronze: exact values (closed form; polar quadrature for the cup). Crosses: the limits $\lambda C(m-1)!$ of \cref{thm:volume:zeta}(ii). Marks: Monte Carlo means of $K^z$ over $10^7$ prior draws (seeds $411$, $412$) with nominal $95\%$ intervals, mean $\pm1.96$ standard errors; left of the dotted line $z=-\frac14$ the variance of $K^z$ is infinite and the intervals mean nothing. Scripts: \texttt{figures/scripts/04-volume-mc.py}, \texttt{04-volume-zeta.py}.}\label{fig:volume:mc}
\end{figure}

\Cref{fig:volume:mc}(a) shows the measured volumes against the exact ones. The exact value lies inside the $95\%$ interval at $74$ of the $75$ pairs of model and level; the miss is the regular plane at $\eps=10^{-5.5}$, with $444$ draws against an expected $497$ (nested counts share their draws, so the next two levels are low as well). The chord slopes between neighbouring levels, with exact $95\%$ intervals from the binomial law of the inner count given the outer one, cover the exact chord slopes at $55$ of the $56$ places for the singular models; the product model's local slope $\frac12-1/(2+\log\frac1{2\eps})$ is $0.331$ at $\eps=10^{-2}$, $0.405$ at $10^{-4}$ and $0.449$ at $10^{-8}$.

Weighted least squares in $\log V$, over the levels $10^{-8}\le\eps\le10^{-2}$ that hold at least $100$ draws ($\eps\ge10^{-6}$ for the regular plane), with $95\%$ intervals from $2000$ parametric-bootstrap replicates of the multinomial counts in the shells between levels (seed $499$), turns the volumes into exponents. A straight line in log--log coordinates gives $1.0006$ $[0.9987,1.0025]$ for the regular plane, $0.5001$ $[0.4990,0.5011]$ for the cup, $0.7494$ $[0.7486,0.7503]$ for the quartic model and $0.6665$ $[0.6653,0.6676]$ for the sextic model. For the product model it gives $0.3905$ $[0.3901,0.3909]$, the average local slope, which understates $\lambda$ by more than a fifth. The leading factor $(\log\frac1\eps)^{m-1}$ alone overcorrects, to $0.520$ $[0.519,0.520]$, because the exact volume $\eps^{1/2}(c_0+c_1\log\frac1\eps)$ has $c_0/c_1=1.31$, not small against $\log\frac1\eps$ on this range. The exact form \eqref{eq:volume:gammatail}, $\eps^\lambda$ times a polynomial of degree $m-1$ in $\log\frac1\eps$, gives $0.5001$ $[0.4988,0.5013]$, with $c_1=0.708$ against the exact $\frac1{\sqrt2}=0.707$.

\Cref{fig:volume:mc}(b) measures the other side of the theorem. The exact curves reach the predicted limits: at $z=-0.4999$ they are at $0.65548$ and $0.35353$, against $0.65551$ and $0.35355$. Plain Monte Carlo of $\zeta(z)=\E[K^z]$ agrees with the exact values to within $0.1\%$ for $z\ge-0.25$ and falls below them further left: at $z=-0.45$ it returns $82\pm15$ for the product model against $136.6$, and $11.2\pm2.1$ for the cup against $12.75$. By parts (i) and~(ii) the variance of $K^z$, $\zeta(2z)-\zeta(z)^2$, is infinite for $z\le-\frac\lambda2$: there the few draws closest to the true set carry most of the integral and are usually missing, so the estimate is typically too small and its nominal interval too narrow. The pole that defines $\lambda$ is invisible to the most direct estimator of $\zeta$.

\section{The ledger}\label{sec:volume:ledger}

The prior mass within $\eps$ of the truth shrinks like $\eps^\lambda(\log\frac1\eps)^{m-1}$, and the two exponents are the learning coefficient and its multiplicity: the exponents of the population free energy (\cref{lem:volume:abelian}), the largest pole of the zeta function and its order (\cref{thm:volume:zeta}), and, for normal crossings, explicit functions of the exponents of the loss and of the prior (\cref{prop:volume:monomial}). The measurements confirm every exponent and show two costs: seeing $\lambda$ by Monte Carlo costs about $\eps^{-\lambda}$ draws, and a multiplicity $m\ge2$ biases every slope by about $(m-1)/\log\frac1\eps$ unless the logarithm is modeled.

\ledgercard{The volume function $V(\eps)=\varphi\{K<\eps\}$, the prior mass of nearly true parameters, and its scaling exponent.}
{$V(\eps)\sim C\eps^{d/2}$: an ellipsoid, $\eps^1$ for two parameters; $\bar F(n)=\frac d2\log n+O(1)$.}
{$V(\eps)\sim C\eps^\lambda(\log\frac1\eps)^{m-1}$ with $-\lambda$ the largest pole of $\zeta$, of order $m$ (\cref{thm:volume:zeta}); $\bar F(n)=\lambda\log n-(m-1)\log\log n+O(1)$ (\cref{lem:volume:abelian}); $\lambda=\min_j\frac{h_j+1}{2k_j}$ for normal crossings (\cref{prop:volume:monomial}).}
{Quartic cup: $V(\eps)=1.311\,\eps^{1/2}$ exactly for $\eps\le\frac12$, so $\lambda=\frac12$, $m=1$, against $\frac d2=1$; $\bar F(n)=\frac12\log n-0.150$. Product model: $V\sim0.707\,\eps^{1/2}\log\frac1\eps$, $\lambda=\frac12$, $m=2$.}
{Monte Carlo, $10^7$ to $10^8$ prior draws: exponents $1.0006$, $0.5001$, $0.5001$ (with the logarithm; $0.39$ without), $0.7494$, $0.6665$ against $1$, $\frac12$, $\frac12$, $\frac34$, $\frac23$; $\zeta$ by Monte Carlo fails for $z<-\lambda/2$.}
{A known population loss $K\ge0$; a volume of the power-logarithm form (guaranteed for analytic $K$ through resolution); a prior with $\varphi\{K=0\}=0$, positive near the true set; bounded $K$ and Delange's Tauberian theorem for the converse.}

\section{Where the theory stops}\label{sec:volume:stops}

\paragraph{What was assumed.} First, that we know $K$. The volume belongs to the population loss, which in practice is unknown; a practitioner has $K_n$, whose sublevel sets fluctuate on the scale $n^{-1/2}$, much larger than the scale $1/n$ at which the free energy reads the volume. For models linear in their coefficients this moves the free energy only by a bounded amount (\cref{thm:freeenergy:linear}); for other models it is part of Watanabe's theorem. Second, the power-logarithm form of $V$: a theorem for analytic $K$; for losses comparable with an analytic one, such as the tanh neuron with Gaussian inputs or the mixture, comparison gives $\lambda$ and $m$ but not $C$; for a ReLU network, whose loss is not analytic, nothing in this chapter guarantees it. Third, a prior that is positive near the true parameters and puts no mass on them. An atom at the truth makes $V(0+)>0$ and the free energy bounded; a prior that vanishes on the true set raises $\lambda$. Jeffreys' prior vanishes wherever the Fisher information is degenerate and turns the quartic cup into the regular plane, with $\lambda=1$ (\cref{ex:volume:jeffreys}), in line with Watanabe's statement that Jeffreys' prior gives $\lambda\ge\frac d2$ \parencite{Watanabe2024}. Fourth, the shape of the parameter set: for the product model with the truth $ab=1$, the true set meets $[-1,1]^2$ only at two corners, $V(\eps)\sim\frac12\eps$, and $\lambda=1$, not the $\frac12$ of a truth whose curve crosses the interior.

\paragraph{Where it breaks.} Volumes cannot be measured where they matter. Plain Monte Carlo needs about $100\,\eps^{-\lambda}/C$ draws for $10\%$ accuracy, $1.3\times10^{13}$ already at $\eps=10^{-2}$ for a regular model with ten parameters, and a network's learning coefficient is a sum of contributions from many directions. The learning coefficients of networks are therefore estimated through the mean energy of a tempered posterior (\cref{ch:temperature,ch:llc}), which the second statement of \cref{thm:toolkit:karamata} ties to the same $\lambda$, and not through volumes; the zeta function is no easier (\cref{fig:volume:mc}(b)). Even exact volumes show their exponents slowly when $m\ge2$, and this qualifies the statement that singular points dominate. \Textcite{Skalse2023} argues, in a blog post, that singular points need not dominate without analyticity and infinite data, with the example $\min(\abs x,\abs y)$. In volume terms the example is sharp: for $K=\frac12\min(a^2,b^2)$ on $[-1,1]^2$, $V(\eps)=2\sqrt{2\eps}-2\eps$, so $\lambda=\frac12$ and $m=1$, and the square $\abs a,\abs b<0.1$ holds a share of $V(\eps)$ that tends to $0.1$, its share of the arms: the centre never dominates. For the analytic product model, with the same $\lambda$, that share tends to one, but slowly: by \cref{comp:volume:log} it is $1-2\log10/(1+\log\frac1{\sqrt{2\eps}})$, which is $0.53$ at $\eps=10^{-8}$ ($0.54$ in a Monte Carlo check with $10^7$ draws) and $0.96$ only at $\eps=10^{-100}$.

\paragraph{Open problems.} Two. First, estimate volumes in high dimension: for analytic $K$ on a box in $\R^d$, give an estimator of $V(\eps)$, or of $\lambda$ through $V$, with controlled relative error at a cost polynomial in $d$ and $\log\frac1\eps$. Nested sampling \parencite{Skilling2006} estimates this kind of function, the prior mass above a sequence of likelihood levels, by drawing from the prior restricted to shrinking sets; near a singular true set those sets are thin and branched, and we found no analysis of the method for singular losses. Second, estimate the multiplicity: give an estimator of $m$ from $N$ prior draws whose error probability decays at a stated rate in $N$, for normal-crossing $K$; we do not know of one. The local estimators of \cref{ch:llc} face the same question at a trained network, where $m$ shifts the estimate of $\lambda$ by about $(m-1)/\log n$.

\begin{exercises}
\item \consolidation (Sums of powers.) (a) Derive \eqref{eq:volume:sums} in detail and check the constants $1.4699$, $1.4457$ and $1.2662$ of \cref{tab:volume:running}. (b) Show that $V(\eps)=\frac\pi6(2\eps)^{3/2}$ for $\frac12(a^2+b^2+c^2)$ on $[-1,1]^3$, and $V(\eps)=1.1107\,\eps^{1/2}$ for $\frac12(a^2+b^2)^2$ on $[-1,1]^2$; explain why the last constant lies between $1.3110/\sqrt2$ and $1.3110$. (c) For the mixed cup $\frac12(a^4+b^6)$, show that $\zeta$ has one pole, simple, at $-\frac5{12}$, with $(z+\frac5{12})\zeta(z)\to0.5276$, and that $\bar F(n)=\frac5{12}\log n-0.1155+o(1)$; compare $\lambda$ with the Newton diagram through $(4,0)$ and $(0,6)$ (\cref{thm:newton:plane}).\label{ex:volume:sums}

\item \consolidation (The Abelian lemma at work.) (a) For the product model, compute the ratio of $Z(n)$ to the prediction of \cref{lem:volume:abelian} at $n=10^4$ and $10^8$ ($1.138$ and $1.069$). (b) Let $K=\frac12(c^2+a^2b^2)$ on $[-1,1]^3$. Show that $Z(n)$ factorizes, that $Z(n)\sim\frac\pi4n^{-1}\log n$, and conclude with Karamata's Tauberian theorem (\cref{thm:toolkit:karamata}) that $V(\eps)\sim\frac\pi4\eps\log\frac1\eps$: $\lambda=1$, $m=2$.\label{ex:volume:abelian}

\item \consolidation (Rules from the volume.) (a) Prove $V_1(\eps/c_2)\le V_2(\eps)\le V_1(\eps/c_1)$ if $c_1K_1\le K_2\le c_2K_1$. (b) For the tanh neuron with $x\sim\Unif[-1,1]$ and the uniform prior on $[-2,2]^2$, $K=\frac{a^2}2\E[\tanh^2(bx)]$; use that $\tanh(u)/u$ decreases on $(0,\infty)$ to show $0.2323\,\frac{a^2b^2}6\le K\le\frac{a^2b^2}6$, and deduce that $V(\eps)$ is within constant factors of $\eps^{1/2}\log\frac1\eps$. (c) Show that if $\varphi$ is continuous and positive at an interior true parameter near which $K$ is of class $C^2$, then $V(\eps)\ge c\,\eps^{d/2}$ for small $\eps$, so $\lambda\le\frac d2$.\label{ex:volume:rules}

\item \consolidation (Monomials.) Let $a=b=1$ in \cref{prop:volume:monomial}, on $[-1,1]^d$. Show that $\zeta(z)=\prod_j2/(2k_jz+h_j+1)$ and $(z+\lambda)^m\zeta(z)\to\lambda C(m-1)!$. Show also that $V(\eps)/\prod_j\frac2{h_j+1}$ is the probability that $\sum_{j:k_j>0}E_j>\log\frac1\eps$ for independent exponential $E_j$ of rates $\lambda_j$.\label{ex:volume:zetamono}

\item \byhand (A prior that vanishes on an arm.) Let $K=\frac12a^2b^4$ on $[-1,1]^2$. (a) With the prior $\frac12\abs a$, show that $\zeta(z)=2^{-z}(z+1)^{-1}(4z+1)^{-1}$ and $V(\eps)=\frac43(2\eps)^{1/4}-\frac13(2\eps)$, and check $C=\frac43\,2^{1/4}$ against \cref{prop:volume:monomial}. (b) With the prior $\frac12\abs b$, show that $\zeta(z)=2^{-z}(2z+1)^{-2}$ and that $V$ is exactly the product model's volume; explain this by the change of variables $c=b\abs b$. A prior that vanishes on the arm $b=0$, where the cusp's $\frac14$ lives, raises $\lambda$ to $\frac12$ and creates a logarithm.\label{ex:volume:arm}

\item \byhand (Three variables.) (a) For $K=\frac12a^2b^4c^6$ on $[-1,1]^3$, find $\lambda=\frac16$, $m=1$ and $C=\frac92\,2^{1/6}$ from \cref{prop:volume:monomial}, and show from the tail of $T_a+2T_b+3T_c$ that $V(\eps)=\frac92s^{1/3}-4s^{1/2}+\frac12s$ with $s=\sqrt{2\eps}$. (b) For $K=\frac12a^2b^4c^4$ find $\lambda=\frac14$, $m=2$ and $C=2^{-3/4}$, and compute $V(\eps)/(\eps^{1/4}\log\frac1\eps)$ at $\eps=10^{-12}$ by quadrature ($0.580$, against $C=0.595$). Why is the convergence so slow?\label{ex:volume:three}

\item \byhand (State densities.) (a) For the product model, check that $\int_0^{1/2}t^zv(t)\dd t=2^{-z}(2z+1)^{-2}$. (b) For the chain of length $L$, show that $\sqrt{2t}\,(L-1)!\,v(t)=(\frac12\log\frac1{2t})^{L-1}$ on $(0,\frac12)$. (c) For a one-coefficient model with coefficient density $\rho$ (\cref{ch:coefficient}) and $K=\frac12(c-c_0)^2$, show that $\sqrt{2t}\,v(t)=\rho(c_0+\sqrt{2t})+\rho(c_0-\sqrt{2t})$.\label{ex:volume:density}

\item \variation (A Gaussian prior.) Let $a,b$ be independent $\Normal(0,1)$ under the prior. Show that the quartic cup has $V(\eps)\sim0.8346\,\eps^{1/2}$ and $\bar F(n)=\frac12\log n+0.302+o(1)$, and, using the density $\frac1\pi K_0(\abs c)$ of $c=ab$ with $K_0(v)=\log\frac2v-\gamma+o(1)$, that the product model has $V(\eps)\sim\frac{\sqrt2}\pi\eps^{1/2}\log\frac1\eps$. Which results of this chapter need a bounded parameter set?\label{ex:volume:gaussian}

\item \variation (Jeffreys' prior.) For the quartic cup, show that the Fisher information at $(a,b)$ is $\diag(4a^2,4b^2)$, so that Jeffreys' prior on $[-1,1]^2$ is $\abs{ab}$; that under it $c=(a^2,b^2)$ is uniform on $[0,1]^2$; and that $V(\eps)=\frac\pi2\eps$ for $\eps\le\frac12$, the volume of the regular plane, so $\lambda=1=\frac d2$. What does this say about the claim that singular models are cheap?\label{ex:volume:jeffreys}

\item \measurement (The cost of Monte Carlo.) (a) With $10^8$ draws in \texttt{04-volume-mc.py}, adapted to $K=\frac12(a^2+b^2+c^2)$ on $[-1,1]^3$, down to which $\eps$ do the counts stay above $100$? (About $7.7\times10^{-5}$.) (b) For $K=\frac12\abs w^2$ on $[-1,1]^d$, compute the draws needed for $10\%$ accuracy at $\eps=10^{-2}$ when $d=10$ and $d=50$ ($1.3\times10^{13}$ and $1.9\times10^{72}$).\label{ex:volume:cost}

\item \measurement (The zeta function near its pole.) (a) Show that $\Var(K^z)=\zeta(2z)-\zeta(z)^2$ is finite exactly for $z>-\lambda/2$. (b) Run \texttt{04-volume-zeta.py} with ten other seeds, record the estimates of $\zeta(-0.45)$ for the product model, compare their median and mean with $136.6$, and explain why an unbiased estimator is usually far too small.\label{ex:volume:zetamc}

\item \research (Volumes in high dimension.) Design an estimator of $\lambda$ for an analytic loss on $[-1,1]^d$ from the volumes of nested sublevel sets, by nested sampling \parencite{Skilling2006} or by Markov chains restricted to $\{K<\eps_k\}$, and analyze its cost on the chain $\frac12(a_1\cdots a_L)^2$ and on rank-one regression (\cref{ch:catalogue}), where $\lambda$ and $m$ are known. The open part is a guarantee at a cost polynomial in $d$ for singular $K$, whose sublevel sets are thin neighbourhoods of branched sets.\label{ex:volume:highdim}

\item \research (Estimating the multiplicity.) For normal-crossing $K$, test $m=1$ against $m=2$ from the empirical volumes of $N$ prior draws at levels $\eps_1>\dots>\eps_k$, using the exact forms \eqref{eq:volume:gammatail}, and find how $N$ must grow to keep both error probabilities below $0.05$ as the levels shrink. Open: the same for general analytic $K$, whose polynomial in $\log\frac1\eps$ has unknown coefficients and whose lower poles add corrections.\label{ex:volume:multiplicity}
\end{exercises}

\notesandrefs
That $\lambda$ is the exponent of a volume runs through Watanabe's work: the review \parencite{Watanabe2024} states that $\lambda$ is the limit of $\log V(\eps)/\log\eps$, and \textcite{WeiEtAl2023}, who argue from the volume that $\lambda$ counts a singular model's effective parameters, cite Theorem~7.1 of \textcite{Watanabe2009} for this characterization. \Textcite{LauEtAl2025} define the local learning coefficient through such volumes near a trained parameter (their Definition~1), as \cref{ch:llc} does.

\Textcite[Lemma~4]{Watanabe2013} gives the leading term of the state density of a normal-crossing function, and the general theory of such expansions is in the second volume of \textcite[pp.~215--262]{ArnoldEtAl1988}; sublevel sets and Newton polyhedra return in \cref{ch:newton}. Th\'eor\`eme~III of \textcite{Delange1954}, a generalization of Ikehara's Tauberian theorem, needs only a nondecreasing function and the behaviour of its Laplace transform on the boundary line of convergence. The proofs of \cref{lem:volume:layercake,lem:volume:abelian}, \cref{prop:volume:monomial} and parts (i)--(iii) of \cref{thm:volume:zeta} are our presentations of standard arguments; the reduction of part~(iii) to Delange's theorem is ours, and we have not seen it in the singular-learning literature.

\paragraph{Translating the literature.} The review \parencite{Watanabe2024} writes $\theta$, $\Theta$ and $\pi(\theta)$ for our $w$, $W$ and $\varphi(w)$, and $v(t)=\int\delta(t-K(\theta))\pi(\theta)\dd\theta$ for the state density, our $v=V'$. \Textcite{LauEtAl2025} write $V(\eps)$ for the Lebesgue volume of $\{w\in B(w^*):L(w)-L(w^*)<\eps\}$, without a prior, and $c\,\eps^{\lambda(w^*)}(-\log\eps)^{m(w^*)-1}$ for its asymptotics; our $V$ with the uniform prior on the ball $B(w^*)$ is theirs divided by its volume. \Textcite{Saito2007} defines the real log canonical threshold of $f$ as the supremum of the $c$ for which $\abs f^{-c}$ is locally integrable; by \cref{thm:volume:zeta}(i) this is our $\sigma_c$, so $\lambda=\operatorname{rlct}(K)$ for a positive prior, and $\lambda=\frac12\operatorname{rlct}(f)$ when $K=f^2$. \Textcite{Lin2011} defines the real log canonical threshold of $K$ through $\int\abs K^{-z}\varphi$, whose poles are at $z=\lambda$ rather than $-\lambda$, and that of an ideal $\langle f_1,\dots,f_k\rangle$ through $(\sum_if_i^2)^{1/2}$, which gives $2\lambda$ when $K$ is comparable with $\sum_if_i^2$. In the language of \textcite{BinghamEtAl1987}, our hypothesis says that $V$ is regularly varying at $0$ with index $\lambda$.

\printchapterbib
\end{refsection}
